\clearpage\section{Relatório consolidado dos ensaios}
O relatório de ensaios de 24/11/2026, com 17 páginas, fornecido pela equipe, acompanha esta versão como documento complementar separado. Para evitar repetição, as evidências individuais permanecem no Anexo A e o PDF original é preservado sem alterações em \texttt{anexos/Relatorio\_Ensaios\_V7.pdf}. O documento inclui o protocolo de referência, procedimentos, resultados, registros T10–T14, fotos da trena e da mochila, capturas seriais, imagens reconstruídas, hashes, ficha consolidada e conclusão. Sua numeração interna e conteúdo foram preservados para rastreabilidade. \href{https://github.com/enzogomes2006/tcd-rf433-esp32/blob/main/relatorio/Relatorio_Ensaios_V7.pdf}{Consultar o relatório original de ensaios no GitHub}.

\subsubsection*{Índice interno do anexo}
\begin{table}[H]\centering\small\caption{Índice do relatório complementar de ensaios}\label{tab:21}
\rowcolors{2}{tcdlight}{white}\begin{tabularx}{\linewidth}{>{\hsize=0.6392\hsize\linewidth=\hsize\raggedright\arraybackslash}X>{\hsize=1.3608\hsize\linewidth=\hsize\raggedright\arraybackslash}X}
\rowcolor{tcdtablehead} \textcolor{black}{\textbf{Páginas do documento de ensaios}} & \textcolor{black}{\textbf{Conteúdo}} \\
1–2 & Objetivos, preparação da bancada e quadro-síntese. \\
3–5 & Procedimentos e resultados T01–T09. \\
6–7 & T10 — distância superior a 1 m, fotografias e logs. \\
8–9 & T11 — obstáculo passivo, comparação e logs. \\
10–11 & T12 — imagem 64\ensuremath{\times}32 A \ensuremath{\to} B e integridade. \\
12–13 & T13 — imagem 128\ensuremath{\times}64 A \ensuremath{\to} B e integridade. \\
14–15 & T14 — imagem 64\ensuremath{\times}32 B \ensuremath{\to} A, diagnóstico e integridade. \\
16–17 & Ficha consolidada, limitações, conclusão e imagem de referência. \\
\bottomrule\end{tabularx}\end{table}
O anexo deve ser interpretado juntamente com o corpo deste relatório, especialmente as ressalvas sobre T06, os tempos estimados, as condições variáveis e as capturas anteriores.

